% ===== 原文 pp.109--112 =====
% [原文 p.109]
\begin{figure}[htbp]
\centering
\includegraphics[width=0.86\textwidth]{fig21.png}
\figtitle{21}{正文中所解释的夹心构造。我们还展示了生成对称性 $\cS$ 的拓扑算符及其以 Y 形联结（junction）的情形；这些算符局域在对称性边界 $\mathfrak{B}^{\mathrm{sym}}_{\cS}$ 上。}
\end{figure}

\para{夹心构造。}SymTFT 的基本思想体现在夹心构造之中。我们先抽象地陈述这一思想，随后再具体地说明它。

考虑一个 $d$ 维理论 $\cT$，它具有某个离散对称性 $\cS$。这可以是一个高阶形式对称性，也可以是一个高阶群对称性（甚至可以是不可逆对称性）。按照夹心构造，我们可以把 $\cT$ 实现为一次区间紧致化，其中包含下列要素：

\begin{itemize}[leftmargin=1.4em,itemsep=1pt,topsep=1pt]
\item \textbf{对称性 TFT}：这是一个 $(d+1)$ 维的 TQFT $\cZ(\cS)$，它只与对称性 $\cS$ 有关，而与理论 $\cT$ 无关。
\item \textbf{对称性边界条件}：这是 SymTFT $\cZ(\cS)$ 的一个拓扑边界条件 $\mathfrak{B}^{\mathrm{sym}}_{\cS}$，它同样只依赖于对称性 $\cS$，而与理论 $\cT$ 无关。
\item \textbf{物理边界条件}：这是 SymTFT $\cZ(\cS)$ 的一个（可能非拓扑的）边界条件 $\mathfrak{B}^{\mathrm{phys}}_{\cT}$，它对不同的理论 $\cT$ 而不同。
\end{itemize}

见图 21。

对称性 $\cS$ 被分离到对称性边界 $\mathfrak{B}^{\mathrm{sym}}_{\cS}$ 上，而理论 $\cT$ 的动力学则生活在边界 $\mathfrak{B}^{\mathrm{phys}}_{\cT}$ 上。

% [Tutor Note: p.109 符号约定：原文 SymTFT 写作花体 $\mathcal{Z}(\mathcal{S})$，PDF 中形似「3(𝒮)」；边界记号是 Fraktur $\mathfrak{B}$，上标 sym/phys 分别指对称性边界与物理边界。核对方法：图 21 左（蓝）面标 $\mathfrak{B}^{\mathrm{sym}}_{\mathcal{S}}$、右（红）面标 $\mathfrak{B}^{\mathrm{phys}}_{\mathcal{T}}$，与正文三个要素一一对应。]

\para{高阶形式对称性的 SymTFT。}设对称性 $\cS$ 是若干离散高阶形式对称性的乘积
\begin{equation}
  \cS=\prod_{i=1}^{n}G^{(p_i)}\,, \tag{4.256}
\end{equation}

% [原文 p.110]
\begin{figure}[htbp]
\centering
\includegraphics[width=0.92\textwidth]{fig22.png}
\figtitle{22}{具有可逆对称性 $\cS$ 的 $d$ 维理论 $\cT$ 的 SymTFT $\cZ(\cS)$ 与物理边界 $\mathfrak{B}^{\mathrm{phys}}_{\cT}$，是通过在 $(d+1)$ 维中规范化（gauging）$\cS$ 而构造出来的。}
\end{figure}

它具有由反常理论
\begin{equation}
  \cI_{d+1}[A_{p_1+1},A_{p_2+1},\cdots,A_{p_n+1}]
  =\exp\left(\ii\int\cL^{\mathrm{eff}}_{d+1}[A_{p_1+1},A_{p_2+1},\cdots,A_{p_n+1}]\right), \tag{4.257}
\end{equation}
所刻画的混合 't Hooft 反常，其中 $A_{p_i+1}$ 是 $G^{(p_i)}$ 对称性的背景场，$\cL^{\mathrm{eff}}_{d+1}$ 是反常理论的有效拉氏量。我们假定 $\cS$ 中出现的任何 0-形式对称性都是阿贝尔的。对于非阿贝尔0-形式对称性，需要对这里讨论的内容作推广。详见 [60]。

于是，夹心构造中涉及 SymTFT 与物理边界的右半部分，可通过规范化所有 $G^{(p_i)}$ 对称性得到。注意，在 $d$ 维中这是做不到的，除非 't Hooft 反常消失；但这里我们是在 $(d+1)$ 维中做规范化，只要把 SPT 相 $\cI_{d+1}$ 附着到理论 $\cT$ 上，这就是可行的。见图 22。

由于一旦忘掉对称性，SPT 相就是平凡理论，规范化过程导致一个离散高阶形式规范理论，但这一理论一般并不属于 2.3 节所讨论的类型。SPT 相的存在把规范理论的作用量修改为
\begin{equation}
  \sum_{i=1}^{n}2\pi\int a_{p_i+1}\cup_\eta\delta b_{d-p_i-1}
  +\int\cL^{\mathrm{eff}}_{d+1}[a_{p_1+1},a_{p_2+1},\cdots,a_{p_n+1}]\,. \tag{4.258}
\end{equation}

现在，为完成夹心构造，我们需要给出对称性边界条件。它由以下两部分组成：

\begin{itemize}[leftmargin=1.4em,itemsep=1pt,topsep=1pt]
\item 对所有规范场 $a_{p_i+1}$ 取 \textbf{Dirichlet} 边界条件（b.c.）；
\item 对所有规范场 $b_{d-p_i-1}$ 取 \textbf{Neumann} 边界条件。
\end{itemize}

% [Tutor Note: 式 (4.258) 中的 $\cup_\eta$ 不能简写成普通 $\cup$：$\eta$ 是 $a_{p_i+1}$ 与 $b_{d-p_i-1}$ 之间的配对（对偶规范群的双线性形式）。量纲/次数核对：$a$ 是 $(p_i+1)$-形式、$\delta b$ 是 $(d-p_i)$-形式，配对后为 $(d+1)$-形式，恰可在 $(d+1)$ 维作用量中积分；若用普通 $\cup$ 则仅在 $\eta$ 为恒等配对时成立。]

% [原文 p.111]
\begin{figure}[htbp]
\centering
\includegraphics[width=0.9\textwidth]{fig23.png}
\figtitle{23}{SymTFT 的体拓扑算符在被推向对称性边界时的命运。}
\end{figure}

的确，此时区间紧致化恰好抵消了对 $G^{(p_i)}$ 对称性的规范化，因为 $\mathfrak{B}^{\mathrm{sym}}_{\cS}$ 对 $a_{p_i+1}$ 取 Dirichlet 边界条件。其结果是，区间紧致化实行替换
\begin{equation}
  a_{p_i+1}\to A_{p_i+1}\,, \tag{4.259}
\end{equation}
并留下 SPT 相 (4.257)。

理论 $\cT$ 的对称性与带电算符都编码在 SymTFT 的拓扑算符之中。回忆 2.3 节的讨论：体内有两种类型的拓扑算符
\begin{equation}
  \begin{aligned}
    U_{\hat g_i}&=\int\hat g_i(a_{p_i+1})\,, &\hat g_i&\in\widehat{G}^{(p_i)}\,,\\
    U_{g_i}&=\int b_{d-p_i-1}(g_i)\,, &g_i&\in G^{(p_i)}\,,
  \end{aligned} \tag{4.260}
\end{equation}
当把它们推到对称性边界上时，Dirichlet 边界条件迫使 $U_{\hat g_i}$ 变得不可见，而由于 Neumann 边界条件，$U_{g_i}$ 得以存活并生成 $\mathfrak{B}^{\mathrm{sym}}_{\cS}$ 的对称性 $\cS$。

% [Tutor Note: (4.260) 的 $\widehat{G}^{(p_i)}$ 是 Pontryagin 对偶群（不是复共轭）；校验：$\hat g_i$ 与 $g_i$ 的配对 $\hat g_i(g_i)$ 在紧致群取值于 $\U(1)$，图 24 第二面板右端正出现因子 $\hat g_i(g_i)\times$。图 23 四联画分别演示 $U_{g_i}$ 被推上边界后存活（上行）与 $U_{\hat g_i}$ 被推上边界后消失（下行，红线化为零）。]

% [原文 p.112]
\begin{figure}[htbp]
\centering
\includegraphics[width=0.72\textwidth]{fig24.png}
\figtitle{24}{沿区间紧致化体拓扑算符 $U_{\hat g_i}$，在理论 $\cT$ 中产生一个算符 $\cO_{\hat g_i}$，它在 $G^{(p_i)}$ $p_i$-形式对称性下具有荷 $\hat g_i\in\widehat{G}^{(p_i)}$。}
\end{figure}

见图 23。

$U_{\hat g_i}$ 在 $\mathfrak{B}^{\mathrm{sym}}_{\cS}$ 上变得不可见，这意味着它们可以终止在 $\mathfrak{B}^{\mathrm{sym}}_{\cS}$ 上。于是我们可把 $(p_i+1)$ 维的 $U_{\hat g_i}$ 算符夹在两个边界之间，从而在理论 $\cT$ 中构造 $p_i$ 维算符；这些算符一般是非拓扑的，因为 $U_{\hat g_i}$ 沿 $\mathfrak{B}^{\mathrm{phys}}_{\cT}$ 的端点一般是非拓扑的。见图 24。它们恰恰就是 $\cT$ 中在 $p_i$-形式对称性下带电的算符！这些算符带电这一事实，来自 $U_{\hat g_i}$ 在体 SymTFT 中在 $U_{g_i}$ 下带电这一事实。见图 24。

% [Tutor Note: p.112 页首的「见图 23。」属于上一页（p.111 末段）的引用，原文因图 24 浮动到本页页首而把它与上文隔开。译文按原文语序保留，不要误以为它指向图 24。]

\para{规范化作为边界条件的改变。}让我们回到一般框架，把拓扑边界条件 $\mathfrak{B}^{\mathrm{sym}}_{\cS}$ 换成 SymTFT $\cZ(\cS)$ 的另一个拓扑边界条件 $\mathfrak{B}^{\mathrm{sym}}_{\cS'}$。边界条件 $\mathfrak{B}^{\mathrm{sym}}_{\cS'}$ 上有拓扑算符生活，它们构成一个对称性 $\cS'$，一般不同于 $\cS$。
